% Vetores unitários em coordenadas esféricas
% Feito para o ÉguaTeX (eguatex.com.br). Uso livre (CC0 1.0).
%
% Compile e pronto: sai só a figura, em PDF.
% Para pôr no seu trabalho, use o PDF com \includegraphics ou copie
% o bloco tikzpicture para dentro de uma figure.
%
% Ângulo theta medido a partir do eixo z; phi no plano xy, a partir do eixo x.
% Os vetores mudam de direção de ponto para ponto: dependem de theta e phi.
\documentclass[tikz,border=4pt]{standalone}
\usepackage[T1]{fontenc}
\usepackage{lmodern}
\usepackage{amsmath}

\colorlet{corVet}{red!75!black}

\begin{document}
\begin{tikzpicture}[
  x={(-0.4435cm,-0.3663cm)}, y={(0.9768cm,-0.1715cm)}, z={(0cm,1cm)},
  line cap=round, line join=round, >=stealth,
  eixo/.style={line width=0.75pt},
  guia/.style={line width=0.55pt, black!35, dashed},
  vetor/.style={->, line width=1.4pt, corVet}]

% Mude estes números e o desenho e os rótulos acompanham.
\def\praio{4} \def\pteta{45} \def\pang{50} % ponto P = (r, theta, phi), ângulos em graus
\def\tam{1.3}                      % tamanho dos vetores, em cm
\pgfmathsetmacro{\st}{sin(\pteta)} \pgfmathsetmacro{\ct}{cos(\pteta)}
\pgfmathsetmacro{\sp}{sin(\pang)} \pgfmathsetmacro{\cp}{cos(\pang)}

% eixos
\draw[eixo] (0,0,0) -- (5,0,0) node[below left] {$x$};
\draw[eixo] (0,0,0) -- (0,5.5,0) node[right] {$y$};
\draw[eixo] (0,0,0) -- (0,0,4.8) node[above] {$z$};

% guias até o ponto
\coordinate (P) at ({\praio*\st*\cp},{\praio*\st*\sp},{\praio*\ct});
\coordinate (Q) at ({\praio*\st*\cp},{\praio*\st*\sp},0);
\draw[guia] (0,0,0) -- (Q) -- (P);
\draw[line width=0.9pt] (0,0,0) -- (P) node[pos=0.6, right=1pt] {$r$};
\draw[->] plot[domain=0:\pteta, samples=30] ({1.2*sin(\x)*\cp},{1.2*sin(\x)*\sp},{1.2*cos(\x)});
\node at ({1.6*sin(\pteta/2)*\cp},{1.6*sin(\pteta/2)*\sp},{1.6*cos(\pteta/2)}) {$\theta$};
\draw[->] plot[domain=0:\pang, samples=30] ({1.3*cos(\x)},{1.3*sin(\x)},0);
\node at ({1.7*cos(\pang/2)},{1.7*sin(\pang/2)},0) {$\phi$};
\fill (P) circle[radius=1.8pt] node[left=3pt] {$P$};

% vetores unitários no ponto P
\draw[vetor] (P) -- ++({\tam*\st*\cp},{\tam*\st*\sp},{\tam*\ct}) node[above right=-2pt] {$\hat{\mathbf{r}}$};
\draw[vetor] (P) -- ++({\tam*\ct*\cp},{\tam*\ct*\sp},{-\tam*\st}) node[below right=-1pt] {$\hat{\boldsymbol{\theta}}$};
\draw[vetor] (P) -- ++({-\tam*\sp},{\tam*\cp},0) node[right] {$\hat{\boldsymbol{\phi}}$};
\end{tikzpicture}
\end{document}
